\section{The instrument, and what it is applied to}
\label{sec:measurement}
\phantomsection\label{sec:design}

\subsection{Five quantities and two floors}
\label{sec:measurement-stack}

A capitulation rate reads only one axis, and no weighted average of the two recovers the
distinction: an average lets a high firmness score offset a low responsiveness one. We
therefore report them separately. \textbf{Installation} asks whether the stance is taken up before
any pressure, and \textbf{hold} what fraction of pressure turns keep it against challenges
that supply no evidence. \textbf{Engagement} asks how a response treats a challenge that does
supply a source, specifics, and a warrant. \textbf{Discrimination} asks whether movement tracks
challenge quality, read as an AUC, and the \textbf{operating point} is the pair of concession
rates at which that AUC is read. Two \textbf{floors} complete the stack: the \emph{adversarial} floor is survival under the
strongest single attack rather than the average one, and the \emph{integrity} floor is the
rate at which a response keeps its position by asserting something false about the challenge.
Both are reported on their own and never averaged into anything else. Installation and hold use a strict rule, counting only
turns at the top of the judges' scale (Appendix~\ref{app:significance}).

The discrimination AUC is the belief/bias distinction as a number: a
quality-blind model scores \(0.5\) \emph{whether it always holds or always folds}, so the
measure cannot be improved by tuning firmness in either direction. A deployment can choose the
operating point; it cannot choose the discrimination. The floors stay outside every aggregate
because an exploit is found once and replayed, and because a model that holds by denying the
record has not held a belief.

\subsection{Interventions and protocol}
\label{sec:design-interventions}
Every intervention conditions the same host (Qwen3.5-4B unless noted) to hold the same assigned
stance; prompt texts are verbatim in Appendix~\ref{app:prompts}. \textbf{Context prompts} at
five strengths (\emph{mild}, \emph{persist}, \emph{engage}, \emph{balanced},
\emph{discerning}, spanning ``argue your position directly'' to explicit never-yield and
explicit update-on-evidence directives; ordered by how hard they push holding, \emph{engage}
and \emph{discerning} invite updating most, \emph{mild} is neutral, and \emph{persist}
forbids yielding). \textbf{Compiled modulation}: a trained write function reads a typed
\emph{belief tree} (claims, arguments, and evidence with credences and SUPPORTS/CONTRADICTS
edges) and emits, in one pass, a fixed rank-16 update to the query and value projections at 11
of the host's 32 layers. The update is installed as a forward pre-hook, applies to every token
of every turn, and is never recomputed mid-conversation, so the belief conditions the forward
pass without occupying the context window. The two projections are built differently: the query
factors come from a content-blind disposition vector, the value factors are compiled from the
host's own hidden states over the belief text, and freezing them one at a time shows the value
channel carries almost all of the hold. Because compilation drops exact statistics and
unfamiliar proper nouns, the deployed configuration pairs the modulation with a short
\emph{curated slot} carrying those specifics in context. The construction sits in the
hypernetwork-adapter family
\citep{chen2025generativeadapter,charakorn2025texttolora,perez1970film}. Everything needed to
reproduce it, the update equations and dimensions, the write function's architecture and
four-stage training curriculum, the belief-tree schema, the exact checkpoints, and the
per-channel ablations, is in Appendices~\ref{app:compiled-internals}--\ref{app:schema} and
\ref{app:qvchannels}.

Arm names compose accordingly: \textsc{ctx}-persist is the persist prompt alone,
\textsc{cmp}+engage is compiled modulation under the engage prompt, and \condF{} is the
deployed configuration, modulation plus slot. \textbf{Matched pairs}, the comparison the paper turns
on: the same prompt, same context, modulation on or off, at every prompt strength.
\textbf{Attempted baselines}: contrastive steering \citep{rimsky2024caa} and a per-belief LoRA;
the steering configuration is degenerate (holds below the unconditioned host) and the LoRA comparison is completed with an objective-matched control in
Section~\ref{sec:objections}.

Protocol: 24 held-out topics, five escalating pressure levels (bare doubt, fabricated
evidence, unsupported flaw assertion, authority, emotional; verbatim templates in
Appendix~\ref{app:l3templates}), three rounds each; pressure turns never restate
the stance. Evidence discrimination uses six fixture types, two well-warranted (real source,
specifics, warrant) and four deficient, style-matched within \(\pm20\%\)
(Appendix~\ref{app:fixtures}); topics are seeded on the side the strongest real evidence
contradicts. Adversarial revocation uses three multi-turn strategies ending in a cold probe
(Appendix~\ref{app:attacks}). Scoring uses a five-judge cross-family ensemble of scriptable models under a rubric that weighs argumentative ground rather than tone (Krippendorff \(\alpha=0.79\)); a sixth judge, reachable only through an agent interface, is used for the headline install and hold cells alone (Appendix~\ref{app:multijudge}). Topic is the unit of analysis (\(n=24\)), intervals are bootstrapped over topics, and paired tests are Wilcoxon with Holm correction. Single-judge cells are marked \(\ddagger\), and an empty generation aborts the run rather than entering a denominator (Appendix~\ref{app:repro}). The canonical results table, with instrument, threshold,
and context configuration stated per cell, is Appendix~\ref{app:canonical}; every number quoted
below is a cell of it or is labelled with its run.
